
\subsection{Data Processing}
The raw data was processed to clean and standardize the transportation network information. 
Key steps included coordinate validation, deduplication, and schema normalization.
Post-cleaning metrics:
\begin{itemize}
    \item Nodes: Unknown
    \item Edges: Unknown
\end{itemize}

\subsection{Network Construction}
The multi-modal transport network is modeled as a directed graph $G=(V,E)$, where $V$ represents intersections and bus stops, and $E$ represents road segments and transfer links.
Edge weights $w_e$ denote travel time in minutes. The graph construction involved:
\begin{enumerate}
    \item Driving Layer: Road segments with speed limits based on highway types.
    \item Transit Layer: Bus stops and routes connected to the driving layer via transfer edges.
\end{enumerate}
Detailed cost functions were applied to ensure realistic routing behavior.

\subsection{Baseline Diagnostics}
Baseline performance was evaluated using $K=100$ randomly sampled OD pairs.
Connectivity and centrality metrics were computed to establish reference performance levels.
Refer to artifacts for detailed distributions.

\subsection{Optimization Model}
We formulate the bus route optimization problem as:
\begin{equation}
    \min f(x) \quad \text{s.t. } x \in X
\end{equation}
where $x$ represents the selection of candidate bus routes.
The objective function includes travel time costs, penalty costs for unreachable pairs, and operational budgets.
A hybrid PSO-GA algorithm was employed to solve this combinatorial optimization problem.
